% Section content template for individual Educative lessons
% This file contains the actual content for a specific section
% Generated from Educative API section content

% Section: Code of Elevator System
% Section ID: 6527052601884672
% Section Slug: code-of-elevator-system
% Generated: 2025-12-11T14:08:48.526660

Using various UML diagrams, we've discussed different aspects of the elevator system and observed the attributes attached to the problem. Let 's explore the more practical side of things, where we will work on implementing the elevator system using multiple languages. This is usually the last step in an object-oriented design interview process.
\\
We have chosen the following languages to write the skeleton code of the different classes present in the elevator control system:

\begin{itemize}
\item
\protect\phantomsection\label{D8UOf5502MmG2mbSxx4C0}
Java
\item
\protect\phantomsection\label{AHP0XCVWWx89UAcDp46lW}
C\#
\item
\protect\phantomsection\label{4LTyluJGJlxFhUKtP3BrE}
Python
\item
\protect\phantomsection\label{64N5Yn2HlbVe_TbIxEAqx}
C++
\item
\protect\phantomsection\label{SmxhQgZ7F4-utY_i6i2xd}
JavaScript
\end{itemize}
\\
If you want to skip directly to the implementation and see the complete workflow, simply scroll down or click \hyperref[Executable-code-Elevator-system]{\textbf{Executable Code: Elevator System}}.

\subsection{Elevator system classes}\label{z-ysLR1kV-tb2lKb3NKiE}
\\
This section will provide the skeleton code of the classes designed in the class diagram lesson.
\begin{quote}
\textbf{Note:} For simplicity, we aren't defining getter and setter functions. The reader can assume that all class attributes are private, accessed through their respective public getter methods, and modified only through their public methods.
\end{quote}
\subsubsection{Enumerations}\label{4BVN4mMZA3hlD0epKFUXP}
\\
First, we will define all the enumerations required in the elevator system. According to the class diagram, three enumerations are used in the system, i.e., \texttt{ElevatorState}, \texttt{Direction} and \texttt{DoorState}\emph{.} The code to implement these enumerations is as follows:
\begin{quote}
\textbf{Note:} JavaScript does not support enumerations, so we will use the Object.freeze()\ method as an alternative that freezes an object and prevents further modifications.
\end{quote}

\textbf{Enum definitions}

\begin{lstlisting}[language=Java]
public enum Direction {
    UP, DOWN, IDLE
}

public enum ElevatorState {
    IDLE, UP, DOWN, MAINTENANCE
}

public enum DoorState {
    OPEN, CLOSED
}

\end{lstlisting}

\subsubsection{Button}\label{RWXA6Mmrdlki4DsobSWvj}
\\
This section contains the implementation of a~\texttt{Button}~class and its subclasses:~\texttt{HallButton}~and \texttt{ElevatorButton}. The ~\texttt{Button}~class has a pure virtual function~\texttt{isPressed()}~in it. The code to implement this relationship is given below:

\textbf{Button and its subclasses}

\begin{lstlisting}[language=Java]
public abstract class Button {
    protected boolean pressed;
    public void pressDown() { }
    public void reset() { }
    public abstract boolean isPressed();
}

public class DoorButton extends Button {
    @Override
    public boolean isPressed() { return pressed; }
}

public class HallButton extends Button {
    private final Direction direction;
    public HallButton(Direction dir) { }
    public Direction getDirection() { return null; }
    @Override
    public boolean isPressed() { return false; }
}

public class ElevatorButton extends Button {
    private final int destinationFloor;
    public ElevatorButton(int floor) { }
    public int getDestinationFloor() { return 0; }
    @Override
    public boolean isPressed() { return false; }
}

public class EmergencyButton extends Button {
    public boolean getPressed() { return false; }
    public void setPressed(boolean val) { }
    @Override
    public boolean isPressed() { return false; }
}
\end{lstlisting}

\subsubsection{Elevator panel and hall panel}\label{tDLcPrXhBdyCmy-YQ43tL}
\\
\texttt{ElevatorPanel} and the \texttt{HallPanel} are classes that use the instance of \texttt{ElevatorButton} and \texttt{HallButton}, respectively.
\\
The code to implement these classes is provided below:

\textbf{The ElevatorPanel and HallPanel classes}

\begin{lstlisting}[language=Java]
import java.util.List;

public class ElevatorPanel {
    private List<ElevatorButton> floorButtons;
    private DoorButton openButton;
    private DoorButton closeButton;
    private EmergencyButton emergencyButton;

    public ElevatorPanel(int numFloors) { }
    public List<ElevatorButton> getFloorButtons() { return null; }
    public DoorButton getOpenButton() { return null; }
    public DoorButton getCloseButton() { return null; }
    public EmergencyButton getEmergencyButton() { return null; }
    public void enterEmergency() { }
    public void exitEmergency() { }
}

public class HallPanel {
    private HallButton up;
    private HallButton down;

    public HallPanel(int floorNumber, int topFloor) { }
    public HallButton getUpButton() { return null; }
    public HallButton getDownButton() { return null; }
}
\end{lstlisting}

\subsubsection{Display}\label{LAMVTngPU8FxxoMdX-5bw}
\\
This component shows the implementation of the \texttt{Display} class. This class is responsible for showing the display inside and outside elevator cars. The code to implement this class is shown below:

\textbf{The Display class}

\begin{lstlisting}[language=Java]
public class Display { private int floor;
private int load; private Direction direction;
private ElevatorState state; private boolean maintenance;
 private boolean overloaded;

 public void update(int f, Direction d, int load, ElevatorState s, boolean overloaded, boolean maintenance) { }
 public void showElevatorDisplay(int carId) { }
}
\end{lstlisting}

\subsubsection{Elevator car}\label{_FM1DBRLxFVNgQAcR3p2f-1}
\\
This section contains the definition of the \texttt{ElevatorCar} class. An elevator car contain instances of \texttt{Door}, \texttt{Display}, and \texttt{ElevatorPanel}. The implementation of this class is represented below:

\textbf{The ElevatorCar class}

\begin{lstlisting}[language=Java]
import java.util.Queue;

public class ElevatorCar {
    private final int id;
    private int currentFloor;
    private ElevatorState state;
    private final Door door = new Door();
    private final Display display = new Display();
    private final ElevatorPanel panel;
    private final Queue<Integer> requestQueue;
    private int load;
    private boolean overloaded;
    private boolean maintenance;

    public ElevatorCar(int id, int numFloors) { }
    public int getId() { return 0; }
    public int getCurrentFloor() { return 0; }
    public ElevatorState getState() { return null; }
    public ElevatorPanel getPanel() { return null; }
    public boolean isInMaintenance() { return false; }
    public boolean isOverloaded() { return false; }
    public void registerRequest(int floor) { }
    public void move() { }
    public void stop() { }
    public void enterMaintenance() { }
    public void exitMaintenance() { }
    public void emergencyStop() { }
    public void addLoad(int kg) { }
    public void removeLoad(int kg) { }
    public Display getDisplay() { return null; }
    public Door getDoor() { return null; }
}

\end{lstlisting}

\subsubsection{Door and floor}\label{SRSE9ivjxp87WdBSivLc7}
\\
This section contains the code for the \texttt{Door} and \texttt{Floor} classes. In the \texttt{Door} class. The enumeration \texttt{DoorState} is used in the \texttt{Door} class and the \texttt{Floor} class contains the instances of \texttt{Display} and \texttt{HallPanel}. The implementation of this class is given below:

\textbf{The Door and Floor classes}

\begin{lstlisting}[language=Java]
import java.util.List;

public class ElevatorSystem {
    private static ElevatorSystem system;
    private final Building building;

    private ElevatorSystem(int floors, int cars) { }
    public static ElevatorSystem getInstance(int floors, int cars) { return null; }
    public List<ElevatorCar> getCars() { return null; }
    public Building getBuilding() { return null; }
    public void callElevator(int floorNum, Direction dir) { }
    public ElevatorCar getNearestIdleCar(int floor) { return null; }
    public void dispatcher() { }
    public void monitoring() { }
}

import java.util.List;

public class Building {
    private final List<Floor> floors;
    private final List<ElevatorCar> cars;

    public Building(int numFloors, int numCars, int numPanels, int numDisplaysPerFloor) { }
    public List<Floor> getFloors() { return null; }
    public List<ElevatorCar> getCars() { return null; }
}


\end{lstlisting}

\subsubsection{Elevator system and building}\label{DpNAhdkSnHg7DcMraAzeH}
\\
The final class of an elevator system is the \texttt{ElevatorSystem} class, which will be a Singleton class, which means that the entire system will have only one instance of this class. Moreover , a \texttt{Building} class contains the instances of \texttt{Floor} and \texttt{ElevatorCar}. The implementation of these Singleton classes is provided below:

\textbf{The ElevatorSystem and Building classes}

\begin{lstlisting}[language=Java]
import java.util. List;

public class ElevatorSystem { private static ElevatorSystem system;
 private final Building building;

 private ElevatorSystem(int floors, int cars) { }
 public static ElevatorSystem getInstance(int floors, int cars) { return null; }
 public List<ElevatorCar> getCars() { return null; }
 public Building getBuilding() { return null; }
 public void callElevator(int floorNum, Direction dir) { }
 public ElevatorCar getNearestIdleCar(int floor) { return null; }
 public void dispatcher() { }
 public void monitoring() { }
}

import java.util. List;

public class Building { private final List<Floor> floors;
 private final List<ElevatorCar> cars;

 public Building(int numFloors, int numCars, int numPanels, int numDisplaysPerFloor) { }
 public List<Floor> getFloors() { return null; }
 public List<ElevatorCar> getCars() { return null; }
}
\end{lstlisting}

\subsection{Executable code: Elevator system}\label{fTyCWVARn6giv_Bc_HbW0}
\\
Below is a fully self-contained, runnable program in Java, C\#, C++, Python, and JavaScript that demonstrates the elevator system's core workflows. The system's main driver code resides in \texttt{Driver.java}, \texttt{Driver.cs}, \texttt{Driver.py}, \texttt{Driver.cpp}, and \texttt{Driver.js} for all respective languages. You can click the ``Run'' button to execute the code.

\subsubsection{What does this code show}\label{A5nDoF6viTeNaQQYFC7sR}

\begin{itemize}
\item
\protect\phantomsection\label{EvOdnOFpswh7z_J4gqOFq}
\textbf{System initialization:} Creates a building with configurable numbers of floors, elevator cars, hall panels, and displays.
\item
\protect\phantomsection\label{_7hKTPPqv1a8zmwTmjN1K}
\textbf{Scenario 1 (Maintenance mode):} Puts an elevator into maintenance, handles a passenger request, and demonstrates how the system dispatches elevators accordingly.
\item
\protect\phantomsection\label{9F0shktEiipnKICVv51lZ}
\textbf{Scenario 2 (Randomized elevator positions):} The system randomly assigns elevator locations and processes a new passenger request, showing how the nearest elevator is assigned and how car positions update.
\item
\protect\phantomsection\label{ndBE5Y7H_sNrebGadJv0E}
\textbf{Status monitoring:} After each action, the system outputs the current status of all elevators and relevant floor displays.
\end{itemize}
\\
\textbf{Hint: Try customizing the code!}
\\
This code is designed for experimentation, not just reading. You can edit, extend, or modify any part of the example.

\textbf{Executable Code: Elevator System}

\begin{lstlisting}[language=Java]
import java.util.Random;
import java.util.List;

public class Driver {
    public static void main(String[] args) {
        int numFloors = 13;
        int numCars = 3;
        int numPanels = 1;    // Number of HallPanels per floor
        int numDisplays = 3;  // Number of Displays per floor

        ElevatorSystem system = ElevatorSystem.getInstance(numFloors, numCars, numPanels, numDisplays);

        // SCENARIO 1
        System.out.println("=== Scenario 1: Elevator 3 in maintenance, passenger calls elevator from floor 7 ===\n");
        system.monitoring();
        System.out.println();

        ElevatorCar car3 = system.getCars().get(2);
        car3.enterMaintenance();
        System.out.println();
        system.monitoring();
        System.out.println();

        runCall(system, 7, Direction.UP);

        car3.exitMaintenance();
        System.out.println("\n--- Resetting maintenance for all elevators ---\n");
        system.monitoring();
        System.out.println();

        System.out.println("=== Scenario 2: Random positions, passenger calls elevator from ground (0) to top (12) ===");

        Random rand = new Random();
        for (ElevatorCar car : system.getCars()) {
            int randomFloor = rand.nextInt(numFloors);
            System.out.printf("\n== Setting random position for Elevator %d ==", car.getId()+1);
            System.out.printf("\n→ Teleporting Elevator %d to floor %d%n", car.getId()+1, randomFloor);
            car.registerRequest(randomFloor);
            car.move();
        }

        System.out.println("\nElevator positions after random repositioning:");
        for (ElevatorCar car : system.getCars()) {
            System.out.printf("Elevator %d ► Floor: %d | State: %s%n",
                              car.getId()+1, car.getCurrentFloor(), car.getState());
        }
        System.out.println();

        runCall(system, 0, Direction.UP);
    }

    private static void runCall(ElevatorSystem system, int floor, Direction dir) {
        System.out.printf("Passenger calls lift on floor %d (%s)%n", floor, dir);
        ElevatorCar nearest = system.getNearestIdleCar(floor);
        if (nearest == null) {
            System.out.println("No idle elevator available right now.");
            return;
        }
        System.out.printf("→ Nearest elevator is %d at floor %d. Lift going %s.%n",
                          nearest.getId()+1, nearest.getCurrentFloor(), dir);
        system.callElevator(floor, dir);
        system.dispatcher();
        System.out.println("\n[Status after dispatch]");
        system.monitoring();
        System.out.println(new String(new char[100]).replace('\0', '-'));
    }
}

\end{lstlisting}

\subsection{Wrapping up}\label{RK4cCJY8yhBHfOE2q80Va}
\\
In this chapter, we've explored the complete design of an elevator control system. We 've also looked at how a basic elevator system can be visualized using various UML diagrams and designed using object-oriented principles and design patterns.

% End of section content